%Version 3.1 December 2024
%%%%%%%%%%%%%%%%%%%%%%%%%%%%%%%%%%%%%%%%%%%%%%%%%%%%%%%%%%%%%%%%%%%%%%
%%                                                                 %%
%% Please do not use \input{...} to include other tex files.       %%
%% Submit your LaTeX manuscript as one .tex document.              %%
%%                                                                 %%
%%%%%%%%%%%%%%%%%%%%%%%%%%%%%%%%%%%%%%%%%%%%%%%%%%%%%%%%%%%%%%%%%%%%%

\documentclass[pdflatex,sn-mathphys-num]{sn-jnl}

%%%% Standard Packages
\usepackage{graphicx}%
\usepackage{multirow}%
\usepackage{amsmath,amssymb,amsfonts}%
\usepackage{amsthm}%
\usepackage{mathrsfs}%
\usepackage[title]{appendix}%
\usepackage{xcolor}%
\usepackage{textcomp}%
\usepackage{manyfoot}%
\usepackage{booktabs}%
\usepackage{algorithm}%
\usepackage{algorithmicx}%
\usepackage{algpseudocode}%
\usepackage{listings}%
\usepackage{url}%

\theoremstyle{thmstyleone}%
\newtheorem{theorem}{Theorem}%
\newtheorem{proposition}[theorem]{Proposition}%

\theoremstyle{thmstyletwo}%
\newtheorem{example}{Example}%
\newtheorem{remark}{Remark}%

\theoremstyle{thmstylethree}%
\newtheorem{definition}{Definition}%

\raggedbottom

%% Listing style for code snippets
\lstset{
  basicstyle=\ttfamily\small,
  breaklines=true,
  frame=single,
  numbers=left,
  numberstyle=\tiny\color{gray},
  keywordstyle=\color{blue},
  commentstyle=\color{green!50!black},
  stringstyle=\color{red!70!black},
  showstringspaces=false,
  tabsize=2
}

\begin{document}

\title[SocGholish Malware Analysis]{Advanced Static Analysis of the SocGholish Malware Family: A Multi-Vector Examination of Fake Browser Update Campaigns}

%%=============================================================%%
%% Author information - placeholder for double-blind review
%%=============================================================%%

\author*[1]{\fnm{Author} \sur{One}}\email{author1@institution.edu}

\author[1]{\fnm{Author} \sur{Two}}\email{author2@institution.edu}

\author[1]{\fnm{Author} \sur{Three}}\email{author3@institution.edu}

\affil*[1]{\orgdiv{Department of Computer Science}, \orgname{Institution Name}, \orgaddress{\city{City}, \state{State}, \country{Country}}}

\abstract{SocGholish is a malware distribution framework that has operated since at least 2018, using compromised websites and fake browser update prompts to push malicious payloads to visitors. The academic literature on it is thin: most available analysis comes from vendor threat reports rather than systematic, peer-reviewed study. We present a static analysis of 160 SocGholish samples from MalwareBazaar, covering JavaScript loaders, PE executables, PowerShell scripts, and HTML lure pages. We built a Python-based analysis framework that extracts over 70 features per sample, including Shannon entropy, obfuscation indicators, embedded network artifacts, and API call patterns. JavaScript loaders show a mean entropy of 5.36, with obfuscation patterns consistent across samples. PE payloads reach entropy values up to 8.0, consistent with packing or encryption. Across the dataset, 54.4\% of samples carried embedded network indicators. The classification engine flagged 11.3\% as suspicious or malicious; the remaining 88.8\% fell below the detection threshold despite being confirmed malware. This outcome reflects SocGholish's evasion design, not a shortcoming of static analysis. These results provide a structured, reproducible dataset and methodology for studying fake browser update campaigns, and we discuss the implications for both detection engineering and threat intelligence.}

\keywords{SocGholish, malware analysis, static analysis, fake browser update, JavaScript obfuscation, Shannon entropy}

\maketitle

%% ============================================================
%% 1. INTRODUCTION
%% ============================================================
\section{Introduction}\label{sec:intro}

Fake browser update campaigns have been a reliable malware distribution method for years, and SocGholish is the most persistent: documented since at least 2018 and still observed as of 2025 \cite{proofpoint2022}. A legitimate website is compromised, JavaScript is injected into its pages, and visitors see a prompt to install a browser update. What actually downloads is a malicious payload, typically a Remote Access Trojan (RAT), a Cobalt Strike loader, or in some campaigns ransomware such as WastedLocker \cite{secureworks2020, mandiant2023}, covering Drive-by Compromise (T1189) and User Execution (T1204) in the MITRE ATT\&CK framework \cite{mitre2024}.

What makes it worth studying closely is the operational discipline behind it, not any novel exploit. Each component is calibrated to the detection environment it faces. The JavaScript is obfuscated just enough to slip past web scanners without looking overtly suspicious; the lure pages pass as genuine browser update dialogs, relying on authority and urgency cues documented in the persuasion literature \cite{cialdini2001, ferreira2015}. Payloads rotate and are sometimes withheld from sandbox-like environments, a practice Miramirkhani et al. \cite{miramirkhani2017} found to be increasingly widespread. This calibrated approach, rather than technical novelty, is what has kept SocGholish running for over seven years. Static, dynamic, and hybrid analysis methods each capture part of the picture \cite{gandotra2014, aslan2020}, yet SocGholish presents a particular challenge because it spans web-based delivery and binary payloads in a way that no single approach handles cleanly.

Pendlebury et al. \cite{pendlebury2019} showed that malware classifiers suffer from temporal bias when trained on narrow snapshots. Characterizing the consistent structural patterns across many samples, obfuscation approach, entropy profile, and embedded network infrastructure, is what makes detection generalize over time. This paper presents that characterization along with a reproducible framework and dataset, and discusses implications for detection engineering and threat intelligence.

This work makes the following contributions:

\begin{itemize}
    \item A reproducible analysis framework for SocGholish samples that handles multiple file types and extracts a broad feature set.
    \item A static analysis of 160 samples characterizing the properties of SocGholish across its delivery chain, from JavaScript loaders to PE payloads.
    \item A comparison of obfuscation techniques and entropy profiles between the JavaScript and executable components.
    \item A discussion of implications for detection engineering and the challenges posed by SocGholish's evasion design.
\end{itemize}

The paper proceeds as follows. Section~\ref{sec:background} reviews the SocGholish framework and related work. Section~\ref{sec:methodology} describes data collection and the analysis framework. Section~\ref{sec:results} reports our findings. Section~\ref{sec:discussion} addresses implications and limitations. Section~\ref{sec:conclusion} concludes.


%% ============================================================
%% 2. BACKGROUND AND RELATED WORK
%% ============================================================
\section{Background and Related Work}\label{sec:background}

\subsection{The SocGholish Framework}\label{subsec:socgholish}

SocGholish, also tracked as FakeUpdates and FAKEUPDATES and attributed to the threat actor TA569, is a distribution system rather than a discrete piece of malware \cite{proofpoint2022}. Figure~\ref{fig:attack_chain} maps its operation against MITRE ATT\&CK stages.

\begin{figure}[h]
\centering
\includegraphics[width=\textwidth]{figures/fig_attack_chain.png}
\caption{SocGholish attack chain mapped to MITRE ATT\&CK framework stages. The upper portion shows the progression from initial compromise through payload delivery, while the lower portion shows the coverage of our analysis framework.}\label{fig:attack_chain}
\end{figure}

Operators target content management systems with known vulnerabilities, injecting a small JavaScript snippet into compromised pages. The snippet is obfuscated and loads an external script from attacker-controlled infrastructure. That script evaluates the visitor's environment, checking browser type, operating system, and geographic location, before deciding whether to serve the fake update prompt \cite{secureworks2020}.

Visitors who pass this check see an overlay modeled after Chrome, Edge, or Firefox update dialogs. The visual fidelity of these pages has improved over time; recent versions are difficult to distinguish from genuine prompts \cite{mandiant2023}. Clicking the button downloads either a JavaScript file or a ZIP archive containing an executable.

Payload selection varies by campaign. NetSupport RAT and Cobalt Strike beacons are the most common; earlier campaigns delivered WastedLocker ransomware \cite{nccgroup2020, cybereason2021}. Using a legitimate remote administration tool like NetSupport makes post-compromise traffic harder to distinguish from normal activity, and combined with code similarity techniques to evade signature-based matching \cite{luo2017}, post-compromise detection becomes genuinely difficult.

\subsection{Related Work}\label{subsec:related}

The methodological foundation for this work sits across several areas of the malware analysis literature. Ye et al. \cite{ye2017} surveyed data mining techniques for malware detection and catalogued the feature representations that had proven most useful up to that point. Souri and Hosseini \cite{souri2018} covered data-driven detection approaches more broadly, mapping how the field moved from signature-based methods toward statistical and learning-based ones. Ucci et al. \cite{ucci2019} focused specifically on machine learning for malware analysis, identifying feature extraction as the step where most classification decisions are effectively made. Taken together, these surveys informed our choice to build a broad, multi-category feature set rather than optimize a single signal.

Static analysis of JavaScript malware has a reasonably long history. Curtsinger et al. \cite{curtsinger2011} built ZOZZLE, a Bayesian classifier over JavaScript AST features that achieved high precision at low computational cost. Rieck et al. \cite{rieck2010} took a behavioral angle, classifying JavaScript based on runtime-observable patterns rather than source structure. Cova et al. \cite{cova2010} focused specifically on drive-by-download attacks, developing methods to detect malicious JavaScript embedded in web pages. Fass et al. \cite{fass2018} introduced JStap, which operates on both syntactic and semantic AST features and works as an efficient pre-filter rather than a standalone detector. The recurring difficulty across all of this work is minimal obfuscation. Likarish et al. \cite{likarish2009} showed that classification on obfuscated JavaScript can be effective, but noted that scripts using only light encoding tend to fall near the decision boundary of syntactic classifiers. Ndichu et al. \cite{ndichu2019} made progress here by embedding AST features as paragraph vectors, capturing distributional patterns that discrete feature counts miss. SocGholish loaders sit precisely in this difficult range: obfuscated enough to resist casual inspection, but not so heavily transformed that they look structurally abnormal.

Entropy has been used as a malware indicator since Lyda and Hamrock \cite{lyda2007} showed that packed and encrypted executables produce measurably higher byte-level entropy than benign software. That observation, grounded in Shannon's information theory \cite{shannon1949}, is now a standard component of most static analysis toolchains. Ugarte-Pedrero et al. \cite{ugarte2015} refined it by showing that section-level entropy within PE files provides a more reliable signal than whole-file values, since packing often affects only specific sections. What the entropy literature has not addressed in depth is the middle range, values between roughly 5.0 and 7.0, where SocGholish's JavaScript loaders consistently fall. Mohaisen et al. \cite{mohaisen2015} made a related observation: the absence of known-bad static indicators does not indicate benignity, particularly for families that deliberately minimize their detectable footprint. This framing is useful for interpreting our classification results. Saxe and Berlin \cite{saxe2015} and Raff et al. \cite{raff2018} pushed toward neural approaches that process raw bytes without hand-crafted features, but both require large labeled corpora that simply do not exist for a targeted family like SocGholish.

The machine learning classification literature offers useful context for the design choices in our framework. Kolter and Maloof \cite{kolter2006} showed that byte-level n-gram features support reasonable classification accuracy on executable malware. Arp et al. \cite{arp2014} developed DREBIN for Android malware and showed that broad, lightweight feature sets can match more complex models while remaining interpretable, a property that matters when results need to be actionable for threat intelligence. Damodaran et al. \cite{damodaran2017}, publishing in this journal, directly compared static, dynamic, and hybrid approaches and found hybrid methods generally outperform either in isolation, a finding that motivates extending our static analysis with dynamic methods in future work. A persistent structural problem is concept drift: Jordaney et al. \cite{jordaney2017} showed that classifiers trained on a snapshot of a malware family degrade as the family evolves, and proposed TRANSCEND to detect when a model's outputs are no longer reliable. For SocGholish, which has been in continuous development since at least 2018, this is not a theoretical concern but a practical one. Stinson and Mitchell \cite{stinson2008} characterized the problem from the attacker's side: operators systematically probe detection boundaries to find the minimum modification needed to evade a classifier. The evasion discipline we observe in SocGholish's loaders is consistent with exactly this kind of iterative refinement.

Research on social engineering and web-based delivery provides additional context. Provos et al. \cite{provos2007} found compromised legitimate websites at the center of drive-by download infrastructure well before SocGholish was documented. Wang et al. \cite{wang2006} built the Strider HoneyMonkey system to automatically patrol exploit-hosting sites, establishing a methodology that anticipates the compromised-site model. The exploit kit literature, surveyed by Nelms et al. \cite{nelms2015}, covers systems that share infrastructure patterns with SocGholish but differ in one critical respect: exploit kits require a vulnerable browser, while SocGholish requires only a willing click. That difference in the trust model is part of why fake update campaigns have proven more durable as browsers patched the vulnerabilities that exploit kits depended on. Kapravelos et al. \cite{kapravelos2013} built Revolver to detect evasive web malware through script similarity, an approach that could complement feature-based analysis like ours. Rahbarinia et al. \cite{rahbarinia2015} tracked malware-control domains across ISP networks, directly relevant to understanding how SocGholish rotates its C2 infrastructure. Thomas et al. \cite{thomas2015} documented deceptive ad injection at scale, with structural parallels to the compromised-site delivery model SocGholish uses.

On obfuscation specifically, Kim et al. \cite{kim2022} found variable renaming, string encoding, and control flow obfuscation to be the most widely used techniques in web-based malware. Skolka et al. \cite{skolka2019} measured baseline obfuscation rates across legitimate web code and found that a substantial fraction of benign scripts use some form of minification or light encoding. That baseline is precisely what SocGholish exploits: its loaders are designed to be statistically indistinguishable from normal minified JavaScript, not to hide through exotic technique. Anderson et al. \cite{anderson2018} showed that reinforcement learning agents can modify PE binaries to evade static classifiers while preserving malicious functionality, and Pierazzi et al. \cite{pierazzi2020} formalized the problem-space constraints that bound what modifications are feasible in practice. Gao et al. \cite{gao2018} demonstrated that contextual signals such as delivery mechanism and domain reputation are often necessary to disambiguate low-signal samples that fall near the detection boundary. Our results point to the same conclusion for SocGholish's JavaScript loaders.


%% ============================================================
%% 3. METHODOLOGY
%% ============================================================
\section{Methodology}\label{sec:methodology}

\subsection{Data Collection}\label{subsec:data_collection}

Samples were drawn from MalwareBazaar \cite{malwarebazaar}, a public repository operated by abuse.ch. We queried five tags associated with SocGholish: \texttt{SocGholish}, \texttt{FakeUpdates}, \texttt{FakeUpdate}, \texttt{FakeBrowserUpdate}, and \texttt{TA569}, along with the \texttt{SocGholish} malware signature. After deduplication by SHA-256 hash, this yielded 196 unique entries. Of these, 160 were successfully downloaded and extracted. The collection spans June 2021 to March 2026.

Table~\ref{tab:dataset} summarizes the composition of our dataset.

\begin{table}[h]
\caption{Dataset composition by file type}\label{tab:dataset}
\vspace{-2mm}
\begin{tabular*}{0.75\textwidth}{@{\extracolsep\fill}lrr@{}}
\toprule
File Type & Count & Percentage \\
\midrule
JavaScript (.js) & 132 & 82.5\% \\
PE Executable (.exe) & 26 & 16.3\% \\
PowerShell (.ps1) & 1 & 0.6\% \\
HTML (.html) & 1 & 0.6\% \\
\midrule
\textbf{Total} & \textbf{160} & \textbf{100\%} \\
\botrule
\end{tabular*}
\vspace{-3mm}
\end{table}

\noindent The skew toward JavaScript reflects how SocGholish operates: every targeted visitor receives a loader, but only a fraction proceed to download the PE payload. JavaScript loaders are therefore far more frequently captured and submitted to public repositories. The PE samples represent second-stage payloads delivered after the initial social engineering stage succeeds.

All samples were handled in an isolated environment on a dedicated remote server. None were executed; all analysis was performed through static methods only.

\subsection{Analysis Framework}\label{subsec:framework}

We built a Python analysis framework that processes each sample through a sequential pipeline. Figure~\ref{fig:architecture} shows the architecture.

\begin{figure}[h]
\centering
\includegraphics[width=\textwidth]{figures/fig_framework_architecture.png}
\caption{Architecture of the SocGholish analysis framework. The system processes samples through four layers: data acquisition, file-type routing, analysis engine, and output generation.}\label{fig:architecture}
\end{figure}

File-type detection runs first, routing each sample to a specialized analyzer. We implemented four:

\begin{itemize}
    \item \textbf{JavaScript Analyzer}: Parses JavaScript source code to extract function counts, variable declarations, obfuscation patterns (eval usage, String.fromCharCode, base64 encoding, hex encoding), embedded URLs, IP addresses, and domain names.
    \item \textbf{PE Executable Analyzer}: Extracts PE header information, section entropy, import table analysis, suspicious API calls (process manipulation, file I/O, registry access, network operations), and packing indicators.
    \item \textbf{PowerShell Analyzer}: Identifies cmdlet usage, encoded command patterns, download cradles, execution policy bypasses, and suspicious PowerShell-specific behaviors.
    \item \textbf{HTML Analyzer}: Extracts embedded scripts, iframe injections, form actions, meta refresh redirects, and social engineering indicators.
\end{itemize}

Every sample, regardless of type, also passes through a common stage that computes:

\subsubsection{Shannon Entropy}\label{subsubsec:entropy}

We compute Shannon entropy $H$ for each sample as a measure of randomness in its byte distribution:

\begin{equation}
H(X) = -\sum_{i=0}^{255} p(x_i) \log_2 p(x_i)
\label{eq:entropy}
\end{equation}

\noindent where $p(x_i)$ is the probability of byte value $x_i$ in the file \cite{shannon1949}. Values range from 0 (all identical bytes) to 8 (perfectly uniform distribution). Plaintext files typically fall between 4.0 and 5.5; packed or encrypted files tend to exceed 7.0 \cite{lyda2007}. Whole-file entropy can miss localized high-entropy regions, so we also examine per-section entropy for PE samples, following Ugarte-Pedrero et al. \cite{ugarte2015}.

\subsubsection{Feature Extraction}\label{subsubsec:features}

The framework extracts over 70 features per sample across the following categories:

\begin{itemize}
    \item \textbf{File metadata}: size, extension, modification timestamp, cryptographic hashes (MD5, SHA-1, SHA-256).
    \item \textbf{Entropy metrics}: whole-file entropy, entropy ratio (entropy/8.0).
    \item \textbf{Obfuscation indicators}: counts of eval/exec calls, base64-encoded strings, hex-encoded strings, Unicode escape sequences, String.fromCharCode invocations.
    \item \textbf{Code structure}: function count, variable count, line count.
    \item \textbf{Network artifacts}: URLs, IP addresses, domain names, HTTP request patterns.
    \item \textbf{API call profile}: categorized counts of process manipulation, file operations, registry access, network, persistence, and cryptographic API calls.
    \item \textbf{SocGholish-specific patterns}: known signature strings, behavioral patterns characteristic of the SocGholish family.
\end{itemize}

\subsubsection{Classification}\label{subsubsec:classification}

Each sample receives a confidence score $s$ computed as a weighted sum of normalized indicator scores:

\begin{equation}
s = \sum_{k=1}^{6} w_k \cdot \hat{f}_k
\label{eq:confidence}
\end{equation}

\noindent where $\hat{f}_k$ represents the normalized value of the $k$-th indicator (clipped to $[0,1]$) and $w_k$ is its corresponding weight. The six indicator components and their weights are:

\begin{itemize}
    \item \textbf{Suspicious API calls} ($w_1 = 0.25$): count of suspicious API calls divided by 10, capped at 1.0.
    \item \textbf{Obfuscation indicators} ($w_2 = 0.20$): count of obfuscation techniques detected, divided by 10, capped at 1.0.
    \item \textbf{Entropy deviation} ($w_3 = 0.20$): for entropy values above 5.0, computed as $(H - 5.0)/3.0$, capped at 1.0.
    \item \textbf{Script commands} ($w_4 = 0.15$): count of scripting-related commands divided by 10, capped at 1.0.
    \item \textbf{Network activity} ($w_5 = 0.10$): count of embedded URLs divided by 20, capped at 1.0.
    \item \textbf{Behavioral patterns} ($w_6 = 0.10$): fraction of active SocGholish-specific behavioral patterns (out of 5 checked patterns).
\end{itemize}

\noindent Weights sum to 1.0 and were assigned based on discriminative power as documented in the malware analysis literature \cite{gandotra2014, ucci2019}. Suspicious API calls carry the highest weight because they are the strongest static indicator of malicious intent; network activity and behavioral patterns carry lower weights because they generate more false positives under static analysis alone.

Based on the resulting confidence score, samples are classified into three categories:

\begin{itemize}
    \item \textbf{Malicious}: confidence score $s \geq 0.7$, indicating strong static indicators of malicious intent.
    \item \textbf{Suspicious}: confidence score $0.3 \leq s < 0.7$, warranting further investigation.
    \item \textbf{Below-threshold}: confidence score $s < 0.3$.
\end{itemize}

\noindent The 0.3 threshold flags samples showing at least two moderately weighted indicators together; the 0.7 threshold requires strong signals across several categories. Both are heuristic. Given a larger labeled dataset, they could be optimized empirically using ROC analysis.

We use ``below-threshold'' rather than ``benign'' deliberately. All 160 samples are confirmed malware; a low score means the sample's static features do not fire our heuristics, not that it is safe. Mohaisen et al. \cite{mohaisen2015} made this point directly: the absence of known-bad indicators is not evidence of benignity, particularly for families that minimize their static footprint by design.

\subsubsection{Analysis Algorithm}\label{subsubsec:algorithm}

Algorithm~\ref{algo:analysis} shows the full procedure applied to each sample. Metadata extraction and entropy computation run first, since both are file-type agnostic and establish baseline features before any format-specific logic executes. The sample is then routed to the appropriate type-specific analyzer: JavaScript, PE, PowerShell, or HTML. After type-specific analysis completes, every sample passes through a common signature matching stage that checks for known SocGholish behavioral patterns, regardless of file type. Feature vector construction and confidence scoring follow. This ordering is intentional: running the common stages last means the signature check can draw on features extracted by the type-specific analyzer, rather than operating on raw file content alone. The pipeline processes a single sample in under two seconds on average, making batch analysis of large collections practical without specialized hardware.

\begin{algorithm}
\caption{SocGholish Sample Analysis Procedure}\label{algo:analysis}
\begin{algorithmic}[1]
\Require Sample file $f$, output directory $O$
\Ensure Feature vector $\mathbf{v}$, classification $c$, confidence score $s$
\State Compute file metadata: size, hashes (MD5, SHA-1, SHA-256)
\State Detect file type $t$ from extension and magic bytes
\State Compute Shannon entropy $H(f)$ using Equation~\ref{eq:entropy}
\If{$t = $ JavaScript}
    \State Extract JS features: functions, variables, eval count
    \State Detect obfuscation: base64, hex, fromCharCode, Unicode
    \State Extract network artifacts: URLs, IPs, domains
\ElsIf{$t = $ PE Executable}
    \State Parse PE headers, extract section entropy values
    \State Analyze import table for suspicious APIs
    \State Check for packing indicators
\ElsIf{$t = $ PowerShell}
    \State Extract cmdlets, encoded commands, download cradles
\ElsIf{$t = $ HTML}
    \State Extract embedded scripts, iframes, form actions
\EndIf
\State Match against SocGholish-specific signature patterns
\State Construct feature vector $\mathbf{v}$ with 70+ features
\State Compute confidence score $s = \sum_{i} w_i \cdot f_i(\mathbf{v})$
\State Classify: $c \leftarrow \begin{cases} \text{malicious} & s \geq 0.7 \\ \text{suspicious} & 0.3 \leq s < 0.7 \\ \text{below-threshold} & s < 0.3 \end{cases}$
\State Export results to JSON, CSV
\end{algorithmic}
\end{algorithm}


%% ============================================================
%% 4. RESULTS
%% ============================================================
\section{Results}\label{sec:results}

The framework processed all 160 samples without errors. Results are organized below by analysis dimension.

\subsection{File Type Distribution and Dataset Characteristics}\label{subsec:file_types}

Figure~\ref{fig:file_type} shows the file type distribution alongside summary metrics. JavaScript files make up 82.5\% of the collection, which follows directly from how SocGholish operates: every targeted visitor receives a loader, so loaders are the most commonly captured and submitted file type. The 26 PE executables (16.3\%) are second-stage payloads, delivered only after a visitor completes the fake update interaction. The single PowerShell script and single HTML lure page round out the dataset, representing less common but operationally real artifacts from SocGholish campaigns.

The imbalance reflects a real-world property of the delivery chain rather than a gap in collection. We report per-type statistics throughout to avoid aggregate figures being dominated by the JavaScript population.

\begin{figure}[h]
\centering
\includegraphics[width=0.95\textwidth]{figures/fig_file_type_distribution.png}
\caption{Distribution of file types in the SocGholish sample dataset (n=160) with key dataset summary metrics.}\label{fig:file_type}
\end{figure}

\subsection{Entropy Analysis}\label{subsec:entropy_results}

Entropy separates the two dominant file types clearly, as Figure~\ref{fig:entropy_kde} shows.

\begin{figure}[h]
\centering
\includegraphics[width=0.9\textwidth]{figures/fig_entropy_kde.png}
\caption{Kernel density estimation of Shannon entropy for JavaScript loaders versus PE executables. The vertical dashed lines indicate the low-obfuscation threshold (5.0) and the high-entropy threshold (7.0). JavaScript samples cluster in the 4.8--6.0 range, while PE executables show a bimodal distribution with peaks around 5.5 and 7.5--8.0.}\label{fig:entropy_kde}
\end{figure}

JavaScript samples showed a mean entropy of 5.36 (SD = 0.18, median = 5.33), ranging from 4.82 to 5.90. This sits above the 4.0--4.8 range typical of unobfuscated JavaScript, pointing to light obfuscation rather than encryption. The tight clustering (IQR = 0.18) indicates a consistent obfuscation approach across the loader population.

PE executables showed a mean of 6.99 (SD = 0.91, median = 6.58), ranging from 4.78 to 8.0. Several samples hit the theoretical maximum of 8.0, a reliable indicator of packing or encryption. The difference between the two populations was statistically significant (Mann-Whitney $U = 143.0$, $z = -7.38$, $p < 0.001$, Cohen's $d = 4.09$). The bimodal PE distribution, with one cluster at moderate entropy and another near 8.0, points to different packaging strategies being used for different payload types.

Figure~\ref{fig:entropy_dist} provides a per-sample view of entropy values across the full dataset.

\begin{figure}[h]
\centering
\includegraphics[width=0.9\textwidth]{figures/fig_entropy_distribution.png}
\caption{Shannon entropy values for all 160 samples, sorted by file type and entropy value. The horizontal reference lines indicate typical entropy thresholds for plaintext (below 5.0) and packed/encrypted content (above 7.0).}\label{fig:entropy_dist}
\end{figure}

\subsection{Obfuscation Analysis}\label{subsec:obfuscation_results}

Figure~\ref{fig:obfuscation_breakdown} breaks down the obfuscation techniques detected across file types.

\begin{figure}[h]
\centering
\includegraphics[width=0.9\textwidth]{figures/fig_obfuscation_breakdown.png}
\caption{Breakdown of obfuscation techniques detected across file types. Base64 encoding is the most prevalent technique in JavaScript samples.}\label{fig:obfuscation_breakdown}
\end{figure}

Base64 encoding was the dominant technique across JavaScript samples, consistent with SocGholish's known practice of encoding strings and code segments to resist signature-based detection. \texttt{eval()} appeared less often than expected, in only a subset of samples, which points to a shift away from eval-based execution toward techniques less likely to raise browser-level warnings. Hex encoding and Unicode escape sequences appeared at lower rates; \texttt{String.fromCharCode} was used sporadically, mostly in samples that combined it with other encoding methods.

The mean obfuscation indicator count per JavaScript sample was 1.23 (SD = 2.31, median = 1.0). PE executables registered zero in all source-level obfuscation categories (mean = 0.0), since they rely on binary packing rather than textual encoding. The difference was statistically significant (Mann-Whitney $U = 689.0$, $z = -4.82$, $p < 0.001$). The low per-sample count is not a measurement artifact; it reflects a deliberate design choice. SocGholish loaders use just enough encoding to avoid basic detection without layering techniques in ways that would look structurally abnormal, exactly what Kim et al. \cite{kim2022} found to be the pattern in effective web malware.

\subsection{Network Indicator Analysis}\label{subsec:network_results}

Embedded network indicators, URLs, IP addresses, and domain names, are among the most directly actionable features for detection. 87 of 160 samples (54.4\%) contained at least one. The mean URL count per sample was 11.44, though this figure is skewed by the PE executables, which tend to embed more strings given the nature of compiled binaries referencing external resources.

The split by file type is telling. 64 of 132 JavaScript loaders (48.5\%) contained embedded URLs, compared to 22 of 26 PE executables (84.6\%). Compiled binaries regularly carry hardcoded C2 addresses, download URLs, and configuration endpoints in their string tables. JavaScript loaders, by contrast, frequently build their URLs at runtime through string concatenation or decoding; the actual C2 destination may never appear as a complete, extractable string in the static file.

Figure~\ref{fig:network_bubble} shows the relationship between entropy, URL count, and file size across all samples.

\begin{figure}[h]
\centering
\includegraphics[width=0.9\textwidth]{figures/fig_network_bubble.png}
\caption{Network indicator analysis. Each point represents a sample, with position determined by entropy (x-axis) and URL count (y-axis). Bubble size corresponds to file size. Color indicates file type.}\label{fig:network_bubble}
\end{figure}

PE executables cluster in the high-entropy, high-URL-count region. JavaScript samples vary more in URL count but stay within a narrower entropy band. The separation between the two populations in this two-dimensional space is clear enough that entropy and network indicator count alone could serve as a lightweight triage step for separating loaders from payloads in a mixed collection.

\subsection{Behavioral Profile Comparison}\label{subsec:radar}

Figure~\ref{fig:radar} shows normalized behavioral profiles across file types on a radar chart.

\begin{figure}[h]
\centering
\includegraphics[width=0.75\textwidth]{figures/fig_radar_comparison.png}
\caption{Radar chart comparing normalized behavioral profiles across file types. Each axis represents a different analysis dimension, normalized to a 0--1 scale.}\label{fig:radar}
\end{figure}

The profiles separate cleanly. JavaScript samples score high on code complexity: a mean of 212.1 functions and 1,367 variable declarations per sample, with a mean line count of 7,789, reflecting the verbosity that obfuscated JavaScript typically produces. PE executables score high on entropy (mean 6.99) but register zero on JavaScript-specific features by definition. The two populations occupy nearly complementary regions of the chart with almost no overlap.

The single PowerShell sample shows elevated script command counts, as expected from cmdlet-heavy syntax. The single HTML sample has low entropy and complexity but carries embedded script elements. Neither file type supports statistical conclusions at n=1, but both confirm that SocGholish campaigns produce a genuine range of artifact types.

\subsection{Classification Results}\label{subsec:classification_results}

The classification engine returned 142 samples (88.8\%) as below-threshold, 17 (10.6\%) as suspicious, and 1 (0.6\%) as malicious. Figure~\ref{fig:classification_flow} shows the full flow from file types through to verdicts.

\begin{figure}[h]
\centering
\includegraphics[width=\textwidth]{figures/fig_classification_flow.png}
\caption{Classification flow diagram showing the progression from input file types through the analysis engine to final verdicts. The detection rate (percentage flagged as suspicious or malicious) was 11.3\%.}\label{fig:classification_flow}
\end{figure}

The mean confidence score across all samples was 0.104 (SD = 0.089). The split by file type tells the real story. JavaScript loaders scored a mean of 0.066 (SD = 0.082, median = 0.042, IQR: 0.031--0.060); PE executables scored a mean of 0.291 (SD = 0.084, median = 0.268, IQR: 0.236--0.330). The difference was statistically significant (Mann-Whitney $U = 146.0$, $z = -7.36$, $p < 0.001$). Only 2 of 132 JavaScript loaders (1.5\%) were flagged, and none were classified as malicious. Among the PE executables, 15 of 26 (57.7\%) were flagged as suspicious. Figure~\ref{fig:violin} shows the per-sample distributions.

\begin{figure}[h]
\centering
\includegraphics[width=\textwidth]{figures/fig_violin_comparison.png}
\caption{Violin plots comparing JavaScript and PE executable samples across three dimensions: Shannon entropy, obfuscation indicator count, and confidence score.}\label{fig:violin}
\end{figure}

The 88.8\% below-threshold rate deserves a direct statement: this is not a failure of the framework. All 160 samples are confirmed malware. A low score means the sample did not fire our static heuristics, which is exactly what SocGholish's JavaScript loaders are built to achieve. They use minimal obfuscation, avoid known-bad API calls, and structure their code to resemble legitimate minified browser scripts. The near-zero flag rate on loaders reflects the selection pressure they face: these are the front-facing components of the campaign, subject to web-based scanning at scale. The PE payloads face a different detection environment and respond accordingly, investing in binary-level packing instead.

\subsection{Feature Correlation Analysis}\label{subsec:correlation}

We examined the correlations between extracted features using a heatmap (Figure~\ref{fig:heatmap}) and a scatter matrix (Figure~\ref{fig:scatter_matrix}).

\begin{figure}[h]
\centering
\includegraphics[width=0.85\textwidth]{figures/fig_correlation_heatmap.png}
\caption{Correlation heatmap of key analysis features. Strong positive correlations are visible between file size and URL count, and between entropy and obfuscation indicators.}\label{fig:heatmap}
\end{figure}

\begin{figure}[h]
\centering
\includegraphics[width=0.85\textwidth]{figures/fig_scatter_matrix.png}
\caption{Scatter matrix showing pairwise relationships between entropy, obfuscation count, URL count, and file size. Diagonal panels show the distribution of each feature by file type.}\label{fig:scatter_matrix}
\end{figure}

Across all 160 samples, the strongest correlation was between entropy and confidence score ($r = 0.633$, $t = 10.29$, $p < 0.001$), as expected given that entropy deviation is a weighted component of the scoring function. File size correlated moderately with confidence score ($r = 0.529$, $t = 7.84$, $p < 0.001$), driven by the larger, more heavily packed PE payloads. Obfuscation indicator count also correlated positively with score ($r = 0.457$, $t = 6.46$, $p < 0.001$).

The negative correlation between entropy and obfuscation indicators ($r = -0.282$, $t = -3.69$, $p < 0.001$) is counterintuitive but has a straightforward explanation. The high-entropy PE executables use binary packing, which our framework does not register as a source-level obfuscation indicator. The lower-entropy JavaScript loaders use encoding techniques that it does detect. The two populations are using fundamentally different evasion methods, and the correlation reflects that split. Entropy also correlated positively with file size ($r = 0.396$, $t = 5.42$, $p < 0.001$), consistent with larger PE payloads being more aggressively packed.

The scatter matrix (Figure~\ref{fig:scatter_matrix}) shows the two populations in largely non-overlapping regions of the feature space. The separation is clearest along the entropy axis. A simple linear classifier trained on entropy alone could separate these populations with high accuracy, which has direct implications for automated triage.

\subsection{Per-Sample Feature Heatmap}\label{subsec:heatmap}

Figure~\ref{fig:detection_heatmap} shows a per-sample feature intensity heatmap for the 50 highest-entropy samples.

\begin{figure}[h]
\centering
\includegraphics[width=\textwidth]{figures/fig_detection_heatmap.png}
\caption{Per-sample detection heatmap showing normalized feature intensities for the top 50 samples by entropy. Brighter colors indicate higher normalized values.}\label{fig:detection_heatmap}
\end{figure}

The high-entropy samples, predominantly PE executables, fire across multiple features at once: high entropy, large file size, and high confidence scores co-occur and form a visually distinct band at the top of the chart. The JavaScript samples tend to show isolated elevation: a sample may score high on obfuscation indicators but low on entropy, or the reverse, but rarely triggers all dimensions together.

The practical distinction is between multi-signal and low-signal targets. PE payloads are multi-signal: multiple features fire simultaneously, making threshold-based detection workable. JavaScript loaders are low-signal: individual features produce weak responses, and no single heuristic reliably catches them. Detecting them with acceptable false positive rates likely requires contextual signals outside the static file itself, such as delivery mechanism, domain reputation, or behavioral distance metrics \cite{gao2018}.

\subsection{File Size Analysis}\label{subsec:size}

Figure~\ref{fig:size_analysis} presents the file size distribution and cumulative density function for JavaScript and PE executable samples.

\begin{figure}[h]
\centering
\includegraphics[width=\textwidth]{figures/fig_size_analysis.png}
\caption{File size analysis on a logarithmic scale. Left: histogram showing the bimodal size distribution between JavaScript loaders and PE payloads. Right: cumulative distribution function (CDF) showing that 90\% of JavaScript samples are smaller than 50 KB, while PE samples range up to several megabytes.}\label{fig:size_analysis}
\end{figure}

JavaScript loaders are consistently small: median 0.9 KB, IQR 0.2--6.0 KB. 90\% fall below 38.7 KB and 95\% below 160 KB. Small files download faster and attract less scrutiny from network appliances. The handful of outliers reaching up to 7.5 MB appear to be concatenated or bundled files that include legitimate library code alongside the malicious payload, a technique that dilutes the malicious signal in both entropy and string analysis.

PE executables cover a much wider range, with a median of 1.8 MB and an IQR of 81 KB to 5.7 MB; the largest sample was 78.7 MB. Smaller PE files tend to be lightweight droppers; larger ones include full RAT installations or bundled components. The two-order-of-magnitude size difference between the two populations is itself a useful first-pass triage feature, consistent with Kolter and Maloof's \cite{kolter2006} observation that file metadata alone can support a lightweight classifier.

Table~\ref{tab:summary} summarizes the key quantitative findings.

\begin{table}[h]
\caption{Summary of key analysis metrics across the full dataset (n=160)}\label{tab:summary}
\begin{tabular*}{\textwidth}{@{\extracolsep\fill}lcccc}
\toprule
Metric & All Samples & JavaScript & PE Executable \\
\midrule
Count & 160 & 132 & 26 \\
Mean Entropy (SD) & 5.625 & 5.36 (0.18) & 6.99 (0.91) \\
Entropy Range & 4.78--8.0 & 4.82--5.90 & 4.78--8.0 \\
Median Entropy & -- & 5.33 & 6.58 \\
Mean Obfuscation Count (SD) & 1.05 & 1.23 (2.31) & 0.00 (0.00) \\
Samples with URLs $>$ 0 & 87 (54.4\%) & 64 (48.5\%) & 22 (84.6\%) \\
Mean Confidence Score (SD) & 0.104 & 0.066 (0.082) & 0.291 (0.084) \\
Classified Suspicious/Malicious & 18 (11.3\%) & 2 (1.5\%) & 15 (57.7\%) \\
\botrule
\end{tabular*}
\end{table}


%% ============================================================
%% 5. DISCUSSION
%% ============================================================
\section{Discussion}\label{sec:discussion}

\subsection{Interpretation of Findings}\label{subsec:interpretation}

The JavaScript loaders sit in an entropy range that is, by design, difficult to classify. A mean of 5.36 is elevated above unobfuscated code but well short of thresholds that flag packing or encryption. The tight clustering (IQR = 0.18) makes clear this is not accidental variation but a consistent operational target. The operators have converged on an obfuscation level that clears web scanners without producing an anomalous file structure, exactly the adversarial feature optimization Anderson et al. \cite{anderson2018} describe in their work on static detection evasion.

The PE payloads operate under different constraints. They do not need to pass web filters; they face endpoint scanning. Binary packing is routine in that environment, which is why entropy values near 8.0 are unremarkable at that stage but would be immediately anomalous in a JavaScript file. The evasion strategy is not uniform across the campaign. It is calibrated per stage to the specific detection mechanisms each component encounters. Pierazzi et al. \cite{pierazzi2020} formalize this as problem-space constraints: the feasible set of adversarial modifications depends on the deployment context, and SocGholish's architecture reflects that constraint precisely.

The 11.3\% detection rate is not a framework failure. It is a measurement. Many JavaScript loaders in our dataset show single-digit VirusTotal detection rates across engines that scan JavaScript, which means the result is a property of the samples, not an artifact of our heuristics. Skolka et al. \cite{skolka2019} found that around 38\% of scripts on the open web use some form of minification or obfuscation. SocGholish loaders are built to stay within that normal range. The operators test iteratively against detection engines, refining loaders until evasion rates are acceptable, a practice documented in industry reports \cite{redcanary2023} and characterized by Stinson and Mitchell \cite{stinson2008} as systematic evaluation of evadability.

Static analysis alone will not reliably catch SocGholish JavaScript loaders. The distinguishing features are too weak in isolation. Effective detection requires contextual signals: how the script was delivered, the reputation of the hosting domain, the behavior of the downloaded payload. Kapravelos et al. \cite{kapravelos2013} showed that structural similarity to known-malicious scripts can surface evasive samples that indicator-based methods miss. For PE payloads, static analysis stays practical; the packing produces entropy signatures that standard tools detect well.

\subsection{Comparison with Prior Work}\label{subsec:comparison}

The entropy gray zone described in Section~\ref{subsec:interpretation} has antecedents in the literature, though JavaScript at the 5.0--5.5 level is largely uncharacterized by prior work. Lyda and Hamrock \cite{lyda2007} documented entropy-based detection of packed executables, but their focus was PE files. Fass et al. \cite{fass2018} are the clearest prior parallel: their JavaScript pre-filtering work identifies this range as precisely the point where entropy alone cannot classify samples and additional syntactic features are required. Our measurements confirm that pattern empirically across 132 JavaScript samples.

The base64-dominant obfuscation profile aligns with Kim et al. \cite{kim2022} and with Ye et al.'s \cite{ye2017} finding that string-level obfuscation remains the most widely deployed evasion technique. Neither study analyzed SocGholish specifically; what we add is a distributional account across a large confirmed-malicious sample set, rather than illustrative examples.

Vendor reports on SocGholish \cite{proofpoint2022, secureworks2020} typically examine individual campaigns. Analyzing 160 samples across four file types makes statistical patterns visible that case study formats cannot surface. The feature set (70+ per sample) is comparable in scope to DREBIN \cite{arp2014}, though oriented toward the web-to-binary delivery chain rather than Android applications.

Damodaran et al. \cite{damodaran2017} showed that hybrid analysis outperforms static-only methods. Our classification results are consistent with that finding: the gap between detected and confirmed samples is the predictable outcome of static features operating against a family engineered to defeat them. Jordaney et al.'s \cite{jordaney2017} concept drift analysis adds a temporal dimension relevant to any operational use of this work. SocGholish has evolved continuously since 2018, and a classifier trained on this dataset will degrade without active monitoring for distributional shift.

\subsection{Limitations}\label{subsec:limitations}

Static analysis has a fixed scope. Without executing samples, dynamic behaviors are outside our view: C2 communication, sandbox evasion, payload deployment logic. The measurements we report describe file-level structure only.

The dataset is 160 samples from MalwareBazaar. That collection reflects what was submitted and tagged by the community; samples that were not submitted, or were submitted without the SocGholish tag, are absent. Generalization beyond this collection requires caution.

Classification thresholds of 0.3 (suspicious) and 0.7 (malicious) were set heuristically based on feature distributions. We did not optimize them through ROC analysis against a held-out labeled set, nor did we run an ablation study to measure each feature class's independent contribution. Results may shift under alternative weighting schemes, and future work should test sensitivity across a range of configurations.

The sample composition is imbalanced: 132 JavaScript samples against 26 PE executables. Per-type statistics reported throughout the paper reduce the distortion, but aggregate metrics such as the overall mean entropy of 5.625 are JavaScript-dominated and should be read accordingly.

Finally, this is a temporal snapshot. The collection window runs June 2021 to March 2026. Jordaney et al. \cite{jordaney2017} and Pendlebury et al. \cite{pendlebury2019} have both shown that static characteristics shift as malware families evolve. The entropy ranges and obfuscation patterns documented here are specific to this window and should be treated as a baseline for that period.

\subsection{Implications for Detection}\label{subsec:detection_implications}

PE detection is the more tractable problem here. Packing that drives entropy toward 8.0 is visible to standard tools, and the 57.7\% PE classification rate reflects that. JavaScript is harder. At mean entropy 5.36 with tight clustering, these loaders fall within the normal range for legitimately minified code. Entropy-based methods will produce high false-negative rates on this component until combined with structural or contextual features.

Obfuscation co-occurrence is a more reliable signal than any individual technique. base64, eval, and fromCharCode each appear routinely in legitimate JavaScript. Finding two or three together in a small file is unusual. Multi-feature classifiers \cite{arp2014} formalize this: detection power comes from the joint distribution of features, not from any threshold applied to a single one.

Network artifacts remain actionable. More than 54\% of samples carried embedded URLs, IPs, or domain names. These indicators change between campaigns, but they integrate well into domain reputation pipelines \cite{rahbarinia2015} for short-term detection and correlation.

File size can reduce the analysis burden at scale. With 90\% of JavaScript loaders under 38.7 KB, size filtering as the first stage of a cascaded pipeline \cite{fass2018} narrows the population requiring deeper analysis without meaningful loss of coverage.

Drift monitoring should be standard for any long-lived adversary. With seven years of operational history, the operators adapt as detection improves. Jordaney et al. \cite{jordaney2017} provide a practical framework for identifying when classifier retraining is necessary; for a family with this history, that mechanism needs to be built in from the start.


%% ============================================================
%% 6. CONCLUSION
%% ============================================================
\section{Conclusion}\label{sec:conclusion}

We analyzed 160 SocGholish samples using a Python-based static analysis framework that extracts over 70 features per sample. The results characterize a family that has remained operationally effective for over seven years partly because its design reflects a practical understanding of where detection works and where it does not.

The central observation from this analysis is one of calibration rather than maximum evasion. JavaScript loaders are built to stay within the entropy range of legitimate minified web content, trading some concealment for durability across diverse scanning environments. PE payloads take the opposite approach, using aggressive packing that prioritizes functional concealment over appearing normal. Each component is matched to the detection environment it faces, and that matching is precise enough to be detectable as a design pattern rather than coincidence. This per-stage calibration, formalized by Pierazzi et al. \cite{pierazzi2020} as problem-space constraints, is the defining structural feature of the SocGholish architecture.

The high proportion of confirmed samples that fell below our static detection threshold is the expected result of this architecture, not a measurement error. It identifies where static features reach their limit against this family, which is necessary information for anyone building detection on top of them.

The dataset and framework support direct extensions. Sandbox execution would expose the runtime behaviors this analysis cannot reach. Supervised classification on the extracted features, with drift monitoring following Jordaney et al. \cite{jordaney2017} and adversarial robustness evaluation following Chen et al. \cite{chen2018}, is the natural next step toward an operational detector. A longitudinal study across the June 2021 to March 2026 collection window would reveal how operator strategies shift in response to detection pressure, addressing the temporal bias concerns raised by Pendlebury et al. \cite{pendlebury2019}.

Fake browser update campaigns persist because they exploit a behavioral pattern that file-level technical controls do not address: users respond to urgency and trust browser branding \cite{cialdini2001, ferreira2015}. The technical analysis here documents the infrastructure beneath that social layer. Understanding what the files look like, how they are structured, and where they conceal their function is the foundation for detection strategies that go beyond individual indicators.


\backmatter

\bmhead{Acknowledgements}

We thank abuse.ch for maintaining the MalwareBazaar platform, which made this research possible by providing open access to malware samples. We also acknowledge the broader threat intelligence community whose sample submissions and tagging enabled our data collection.

\section*{Declarations}

\begin{itemize}
\item \textbf{Funding}: This research received no external funding.
\item \textbf{Conflict of interest}: The authors declare no conflict of interest.
\item \textbf{Ethics approval}: Not applicable. This study analyzes publicly available malware samples and does not involve human subjects.
\item \textbf{Data availability}: The SHA-256 hashes of all analyzed samples are available from MalwareBazaar (https://bazaar.abuse.ch). The analysis results and extracted features will be made available upon request.
\item \textbf{Code availability}: The analysis framework source code will be made available upon acceptance of this paper.
\item \textbf{Author contribution}: All authors contributed to the design of the study, data collection, analysis, and manuscript preparation.
\end{itemize}

\begin{appendices}

\section{Sample Code: Entropy Calculation}\label{secA1}

The following Python function implements the Shannon entropy calculation used in our analysis framework:

\begin{lstlisting}[language=Python, caption={Shannon entropy calculation function}]
import math
from collections import Counter

def calculate_entropy(data):
    """Calculate Shannon entropy of byte data."""
    if not data:
        return 0.0

    byte_counts = Counter(data)
    total = len(data)
    entropy = 0.0

    for count in byte_counts.values():
        if count > 0:
            probability = count / total
            entropy -= probability * math.log2(probability)

    return entropy
\end{lstlisting}

\section{Sample Code: Obfuscation Detection}\label{secA2}

The following code excerpt shows the obfuscation detection logic for JavaScript samples:

\begin{lstlisting}[language=Python, caption={JavaScript obfuscation detection}]
import re

def detect_obfuscation(content):
    """Detect common obfuscation techniques."""
    indicators = 0
    details = []

    # Check for eval usage
    eval_count = len(re.findall(
        r'\beval\s*\(', content))
    if eval_count > 0:
        indicators += 1
        details.append(f"eval: {eval_count}")

    # Check for String.fromCharCode
    fcc = len(re.findall(
        r'String\.fromCharCode', content))
    if fcc > 0:
        indicators += 1
        details.append(f"fromCharCode: {fcc}")

    # Check for base64 patterns
    b64 = re.findall(
        r'[A-Za-z0-9+/]{40,}={0,2}', content)
    if len(b64) > 2:
        indicators += 1
        details.append(f"base64 strings: {len(b64)}")

    # Check for hex encoding
    hex_strings = re.findall(
        r'\\x[0-9a-fA-F]{2}', content)
    if len(hex_strings) > 10:
        indicators += 1
        details.append(f"hex escapes: {len(hex_strings)}")

    return indicators, details
\end{lstlisting}

\end{appendices}

\bibliography{sn-bibliography}

\end{document}
